\documentclass[12pt, a4paper, oneside]{ctexart}
\usepackage{amsmath, textcomp,siunitx,amsthm, amssymb, bm, xcolor, tikz, framed, graphicx, hyperref, mathrsfs}

\title{Statistical Mechanics Homework (Chapter5)}
\author{Tang Yi \\ [学号已脱敏]}
\usepackage[a4paper, left=2cm, right=2cm, top=2cm, bottom=2cm]{geometry}
\linespread{1.3} % 行距更紧凑
\definecolor{shadecolor}{RGB}{241, 241, 255}

\newcounter{problemname}
\newenvironment{problem}
  {\begin{shaded}\stepcounter{problemname}\par\noindent\textbf{题目 \theproblemname. }}
  {\end{shaded}\par}
\newenvironment{solution}
  {\par\noindent\textbf{题目 \theproblemname 的解答. }}
  {\par}
\newenvironment{note}
  {\par\noindent\textbf{题目 \theproblemname 的注记. }}
  {\par}


\newcommand{\dbar}{\; \bar{} \hspace{-0.3em} \mathrm d}
\newcommand{\tpar}[3]{\left(\dfrac{\partial #1}{\partial #2}\right)_{#3}}
\begin{document}

\maketitle

\begin{problem}{5-1}
  一维谐振子的能谱为$\varepsilon_n = \left( n+\dfrac{1}{2} \right) h\nu , n = 0,1,2,\cdots$，而且系统的温度足够低。

  (a)求振子处于第一激发态与基态的概率之比；

  (b)如果振子仅可能占据基态和第一激发态，计算平均能量。
\end{problem}

\begin{solution}
  设谐振子的能谱为
\[
\varepsilon_n=\left(n+\frac12\right)h\nu ,\qquad n=0,1,2,\ldots
\]
温度足够低，令 \(\beta=\dfrac{1}{k_B T}\)。

(a) 第一激发态与基态的概率之比为
\[
\frac{P_1}{P_0}
=\exp\!\left[-\beta(\varepsilon_1-\varepsilon_0)\right]
=\exp(-\beta h\nu).
\]

(b) 若系统只可能处于 \(n=0,1\) 两个能级，配分函数为
\[
Z=e^{-\beta\varepsilon_0}+e^{-\beta\varepsilon_1}
=e^{-\beta\varepsilon_0}\left(1+e^{-\beta h\nu}\right).
\]
平均能量为
\[
\langle \varepsilon\rangle
=\frac{\varepsilon_0 e^{-\beta\varepsilon_0}
+\varepsilon_1 e^{-\beta\varepsilon_1}}{Z}
=\varepsilon_0+h\nu\,\frac{e^{-\beta h\nu}}
{1+e^{-\beta h\nu}}.
\]
代入 \(\varepsilon_0=\frac12 h\nu\)，得
\[
\langle\varepsilon\rangle
=\frac12 h\nu+\frac{h\nu}{e^{\beta h\nu}+1}.
\]
\end{solution}

\begin{problem}{5-2}
  服从玻耳兹曼分布的某种理想气体粒子的能量和动量关系为$\varepsilon = cp$，其中$c$为光速。若单位体积中包含$ N $个粒子，计算此气体的定容热容。
\end{problem}

\begin{solution}
  设单粒子能量与动量关系为 \(\varepsilon=cp\)，系统服从玻尔兹曼分布，单位体积中有 \(N\) 个粒子。取 \(\beta=1/(k_B T)\)。

单粒子的配分函数（单位体积，经典近似）为
\[
Z_1=\frac{1}{h^3}\int_{\mathbb{R}^3} e^{-\beta c p}\,d^3p
=\frac{4\pi}{h^3}\int_0^\infty p^2 e^{-\beta c p}\,dp.
\]
利用积分公式 \(\displaystyle\int_0^\infty x^2 e^{-ax}\,dx=\frac{2}{a^3}\)（\(a>0\)），得
\[
Z_1=\frac{4\pi}{h^3}\cdot\frac{2}{(\beta c)^3}
=\frac{8\pi}{h^3}\frac{1}{(\beta c)^3}.
\]

总体配分函数为 \(Z=Z_1^N/N!\)（经典近似），系统总能量为
\[
U=-\frac{\partial}{\partial\beta}\ln Z
=-N\frac{\partial}{\partial\beta}\ln Z_1.
\]
由 \(\ln Z_1=\ln\!\big(\tfrac{8\pi}{h^3 c^3}\big)-3\ln\beta\)，得到
\[
U=-N\left(-\frac{3}{\beta}\right)=\frac{3N}{\beta}=3N k_B T.
\]

因此定容热容为
\[
C_V=\left(\frac{\partial U}{\partial T}\right)_V=3N k_B.
\]

\end{solution}

\begin{problem}{5-3}
  有$ N $个质量为$ m $的单原子分子组成的理想气体，它们被装在边长为$ L $的立方容器内，其上下与地底表面平行，设在容器的范围内重力场均匀。气体处于温度为$ T $的平衡态，服从玻耳兹曼分布。试计算：

(a) 分子的平均动能； (b) 分子的平均势能。
\end{problem}

\begin{solution}
  设有 \(N\) 个质量为 \(m\) 的单原子分子置于边长为 \(L\) 的立方容器内，重力加速度为 \(g\)，取竖直方向 \(z\) 从底面（\(z=0\)）向上到顶面（\(z=L\)）。系统处于温度 \(T\) 的热平衡，服从玻尔兹曼分布，记 \(\beta=1/(k_B T)\)。

单个分子的配分函数可分为动量部分与空间部分的乘积：
\[
Z_1=\frac{1}{h^3}\int_{\mathbb{R}^3} e^{-\beta\frac{p^2}{2m}}\,d^3p
\int_{V} e^{-\beta m g z}\,d^3r.
\]
动量积分为高斯积分：
\[
\int_{\mathbb{R}^3} e^{-\beta\frac{p^2}{2m}}\,d^3p
=\left(\int_{-\infty}^{\infty} e^{-\beta\frac{p_x^2}{2m}}dp_x\right)^3
=\left(\sqrt{\frac{2\pi m}{\beta}}\right)^3
=\left(\frac{2\pi m}{\beta}\right)^{3/2}.
\]
空间积分因为横向面积为 \(L^2\)，在 \(z\) 上为指数衰减：
\[
\int_{V} e^{-\beta m g z}\,d^3r
= L^2\int_0^L e^{-\beta m g z}\,dz
= L^2\frac{1-e^{-\beta m g L}}{\beta m g}.
\]
因此
\[
Z_1=\frac{1}{h^3}\left(\frac{2\pi m}{\beta}\right)^{3/2}
L^2\frac{1-e^{-\beta m g L}}{\beta m g}.
\]

(a) 平均动能。动能只来自动量部分，平均动能可由动量配分函数给出：
\[
\langle K\rangle
=-\frac{\partial}{\partial\beta}\ln\Big[\,(2\pi m/\beta)^{3/2}\,\Big]
=\frac{3}{2\beta}=\frac{3}{2}k_B T.
\]

(b) 平均势能。只看位置的配分函数
\[
Z_{\text{pos}}=L^2\frac{1-e^{-\beta m g L}}{\beta m g},
\]
单粒子的平均势能为
\[
\langle V\rangle
=-\frac{\partial}{\partial\beta}\ln Z_{\text{pos}}
=\frac{1}{\beta}-\frac{m g L}{e^{\beta m g L}-1}.
\]

\end{solution}

\begin{problem}{5-4}
  1mol 带电荷为$ q $的单原子分子组成的理想气体，在边长$ L $的立方体中运动，其中心沿$ x $方向存在一个均匀电场$ E$，分子的能量为 $\varepsilon = \dfrac{p^2}{2m} - Eqx$。求出该气体的 (a) 配分函数；(b) 内能。
\end{problem}

\begin{solution}
    设气体为 1 mol（粒子数为 \(N=N_A\)），单粒子能量为
\[
\varepsilon=\frac{p^2}{2m}-E q x,
\]
系统处于温度 \(T\)，记 \(\beta=\frac{1}{k_B T}\)。单粒子配分函数可分为动量和位置两部分：
\[
Z_1=\frac{1}{h^3}\int_{\mathbb{R}^3}e^{-\beta\frac{p^2}{2m}}\,d^3p
\int_{V}e^{\beta E q x}\,d^3r.
\]
动量积分为高斯积分：
\[
\int_{\mathbb{R}^3}e^{-\beta\frac{p^2}{2m}}\,d^3p
=\left(\frac{2\pi m}{\beta}\right)^{3/2}.
\]
空间积分在立方体 \(0\le x,y,z\le L\) 上为
\[
\int_{V}e^{\beta E q x}\,d^3r
= L^2\int_0^L e^{\beta E q x}\,dx
= L^2\frac{e^{\beta E q L}-1}{\beta E q}.
\]
因此单粒子配分函数为
\[
Z_1=\frac{1}{h^3}\left(\frac{2\pi m}{\beta}\right)^{3/2}
L^2\frac{e^{\beta E q L}-1}{\beta E q}.
\]

(a) 全体系（经典近似）的配分函数为正则配分函数（可交换粒子）：
\[
Z=\frac{Z_1^{\,N}}{N!},\qquad N=N_A\ \text{（1 mol）}.
\]

(b) 内能由
\[
U=-\frac{\partial}{\partial\beta}\ln Z
=-N\frac{\partial}{\partial\beta}\ln Z_1
\]
给出。写出 \(\ln Z_1\)：
\[
\ln Z_1=\frac{3}{2}\ln\!\left(\frac{2\pi m}{\beta}\right)
+\ln L^2+\ln\!\big(e^{\beta E q L}-1\big)-\ln(\beta E q)-3\ln h.
\]
对 \(\beta\) 求导得
\[
\frac{\partial}{\partial\beta}\ln Z_1
=-\frac{3}{2}\frac{1}{\beta}
+\frac{E q L\,e^{\beta E q L}}{e^{\beta E q L}-1}-\frac{1}{\beta}.
\]
于是
\[
U=-N\left(-\frac{5}{2}\frac{1}{\beta}
+\frac{E q L\,e^{\beta E q L}}{e^{\beta E q L}-1}\right)
= N\left(\frac{5}{2}k_B T-\;E q L\;\frac{e^{\beta E q L}}{e^{\beta E q L}-1}\right).
\]

其中，$N = N_A$。
\end{solution}

\begin{problem}{5-5}
  气柱的高度为 H，处在重力场中。试证明此气柱的内能和热容为

  \centerline{$U = U_0 + N k_B T - \frac{N m g H}{e^{\frac{m g H}{k_B T}} - 1}$}

  \centerline{$C_V = C_V^0 + N k_B - \frac{N (mgH)^2 e^{\frac{mgH}{k_B T}}}{\left( e^{\frac{mgH}{k_B T}} - 1 \right)^2} \frac{1}{k_B T^2}$}
\end{problem}

\begin{solution}
  设单粒子能量为
\[
\varepsilon=\frac{p_x^2+p_y^2+p_z^2}{2m}+mg z,\qquad 0\le z\le H.
\]
配分函数可分解为动量部分与位置部分：
\[
Z_1=Z_p\,Z_z,
\]
其中动量积分给出
\[
Z_p=\frac{1}{h^3}\!\left(\frac{2\pi m}{\beta}\right)^{3/2},
\]
竖直方向的重力势能部分
\[
Z_z=\int_0^H e^{-\beta mg z}\,dz
=\frac{1-e^{-\beta mg H}}{\beta mg}.
\]
因此单粒子配分函数为
\[
Z_1=Z_{\rm int}\;\frac{1}{h^3}\!\left(\frac{2\pi m}{\beta}\right)^{3/2}
\frac{1-e^{-\beta mg H}}{\beta mg},
\]
其中 \(Z_{\rm int}\) 为内部自由度配分函数，对应能量记为 \(U_0\)。

系统总配分函数为 \(Z=Z_1^{\,N}/N!\)。  
内能：
\[
U=-\frac{\partial}{\partial\beta}\ln Z
=-N\frac{\partial}{\partial\beta}\ln Z_1.
\]

利用
\[
\ln Z_1=\ln Z_{\rm int}+\frac{3}{2}\ln\!\left(\frac{1}{\beta}\right)
+\ln(1-e^{-\beta mgH})-\ln(\beta mg),
\]
对 \(\beta\) 求导，并用 \(1/\beta=k_B T\)，得
\[
U
=U_0+N k_B T - \frac{N mgH}{e^{\beta mgH}-1}.
\]

即
\[
\boxed{
U=U_0+N k_B T-\frac{N mgH}{e^{\frac{mgH}{k_B T}}-1}
}.
\]

再对 \(U\) 对 \(T\) 求导，注意
\[
x=\frac{mgH}{k_B T},
\quad
\frac{d}{dT}\frac{1}{e^x-1}
=-\frac{(mgH)e^x}{k_B T^2 (e^x-1)^2},
\]
得到
\[
\boxed{
C_V=C_V^0
+N k_B
-\frac{N(mgH)^2e^{\frac{mgH}{k_B T}}}{(e^{\frac{mgH}{k_B T}}-1)^2}
\frac{1}{k_B T^2}
}.
\]
\end{solution}

\begin{problem}{5-6}
  有 N 个全同的、近独立二维简谐振子组成的系统，服从玻耳兹曼统计，已知每 个振子的能谱为

  \centerline{$\varepsilon_{n_1, n_2} = \left(n_1 + \dfrac{1}{2}\right) h \nu_1 + \left(n_2 + \dfrac{1}{2}\right) h \nu_2, \quad n_1, n_2 = 0, 1, 2, \ldots$}

  试求一个振子的平均能量，并给出当温度满足条件$h\nu_1 \ll k_B T \ll h\nu_2$时，它的近似表达式。
\end{problem}

\begin{solution}
  设 \(\beta=\dfrac{1}{k_B T}\)。单个二维简谐振子的能级为
\[
\varepsilon_{n_1,n_2}
=\left(n_1+\frac12\right)h\nu_1+\left(n_2+\frac12\right)h\nu_2,
\qquad n_1,n_2=0,1,2,\ldots
\]
配分函数因两个方向独立而可写成乘积：
\[
Z_1=Z(\nu_1)Z(\nu_2),\qquad
Z(\nu)=\sum_{n=0}^\infty e^{-\beta h\nu(n+1/2)}
=\frac{e^{-\beta h\nu/2}}{1-e^{-\beta h\nu}}.
\]
单振子的平均能量为
\[
\langle E\rangle
=-\frac{\partial}{\partial\beta}\ln Z_1
=h\nu_1\Big(\tfrac12+\frac{1}{e^{\beta h\nu_1}-1}\Big)
+ h\nu_2\Big(\tfrac12+\frac{1}{e^{\beta h\nu_2}-1}\Big).
\]
写成更紧凑的形式：
\[
\boxed{\;
\langle E\rangle
= h\nu_1\left[\frac{1}{2}+\frac{1}{e^{\beta h\nu_1}-1}\right]
+ h\nu_2\left[\frac{1}{2}+\frac{1}{e^{\beta h\nu_2}-1}\right]\; }.
\]

当温度满足 \(h\nu_1\ll k_B T\ll h\nu_2\)时，可作近似：
\[
\frac{1}{e^{\beta h\nu_1}-1}\approx\frac{1}{\beta h\nu_1},\qquad
\frac{1}{e^{\beta h\nu_2}-1}\approx e^{-\beta h\nu_2}.
\]
于是单振子的平均能量近似为
\[
\boxed{\;
\langle E\rangle \simeq k_B T + \frac{1}{2}h\nu_2 + h\nu_2\,e^{-\beta h\nu_2}\; }.
\]
\end{solution}

\begin{problem}{5-7}
  经典转子的能量为

  \centerline{$\epsilon = \dfrac{1}{2 I} p_{\theta}^2 + \dfrac{1}{2 I \sin^2 \theta} p_{\phi}^2$}

  计算相体积、配分函数和定容热容。
\end{problem}

\begin{solution}
  设经典线性转子（质点间距极小、仅有转动自由度）转动惯量为 \(I\)，广义角坐标为 \((\theta,\phi)\)，共轭动量为 \((p_\theta,p_\phi)\)，哈密顿量为
\[
\epsilon=\frac{p_\theta^2}{2I}+\frac{p_\phi^2}{2I\sin^2\theta},\qquad
0\le\theta\le\pi,\;0\le\phi<2\pi.
\]

一、相空间体积（能量小于等于 \(E\) 的相空间体积）
\[
\Gamma(E)=\int_{H\le E} d\theta\,d\phi\,dp_\theta\,dp_\phi.
\]
对于固定的 \(\theta\)，动量空间上满足
\[
\frac{p_\theta^2}{2I}+\frac{p_\phi^2}{2I\sin^2\theta}\le E
\]
的区域是一个半轴长为 \(\sqrt{2I E}\) 与 \(\sqrt{2I E}\sin\theta\) 的椭圆，其面积为
\(\pi\cdot( \sqrt{2IE})(\sqrt{2IE}\sin\theta)=2\pi I E\sin\theta\)。
故
\[
\Gamma(E)=\int_0^{2\pi}d\phi\int_0^\pi d\theta\; \big(2\pi I E\sin\theta\big)
=2\pi I E\cdot 2\pi\int_0^\pi \sin\theta\,d\theta
=8\pi^2 I E.
\]
于是能量密度（对能量求导）为
\[
\frac{d\Gamma}{dE}=8\pi^2 I.
\]

二、单个转子的配分函数（经典近似）
\[
Z_1=\frac{1}{h^2}\int e^{-\beta\epsilon}\,d\theta\,d\phi\,dp_\theta\,dp_\phi
=\frac{1}{h^2}\int_0^\infty e^{-\beta E}\frac{d\Gamma}{dE}\,dE
=\frac{1}{h^2}\,8\pi^2 I\int_0^\infty e^{-\beta E}dE
=\frac{8\pi^2 I}{\beta h^2}.
\]
记 \(\beta=1/(k_B T)\)，可写为
\[
\boxed{\,Z_1=\dfrac{8\pi^2 I k_B T}{h^2}\, }.
\]

若系统有 \(N\) 个互不相互作用的全同经典转子，则总体配分函数（经典近似，可交换粒子）为
\[
Z=\frac{Z_1^N}{N!}.
\]

三、定容内能与定容热容
单个转子的平均能量
\[
\langle\epsilon\rangle
=-\frac{\partial}{\partial\beta}\ln Z_1
=\frac{1}{\beta}=k_B T.
\]
因此总内能与定容热容为
\[
\boxed{\,U=N k_B T,\qquad C_V=\left(\frac{\partial U}{\partial T}\right)_V
= N k_B\, }.
\]

说明：上述结果与等分配定理一致——线性转子有两个独立的转动自由度，每个自由度贡献 \(\tfrac12k_B T\)，因而每个分子转动部分平均能量为 \(k_B T\)。
\end{solution}
\end{document}